\section{Audited claim map and the next arithmetic obligations}\label{sec:register}
\subsection{What this continuation adds to the progress paper}
The conductor branch now has a single normalization map, a proof of the necessary period divisibility, an expanded source self-pair calculation, and a precise conditional rank-nine filter. The actual-order certificate remains visible. The full and skipped rank-eight profiles show why even several correct conductor invariants cannot establish the needed native quotient.

The classical induction branch now identifies its field tower, tower-independence input, tame descent and central normalization separately. Its characteristic extension is named, and the failed fresh access to the underlying positive-characteristic source is not concealed. This leaves a meaningful selected input audit without overstating it as a new proof of all local correspondence results.

The marking branch now contains the full center-fixing extension criterion, the dihedral/quaternion distinction, the sharp free-embedding capacity test, and the exact rank-eight Arf condition. The explicit transvection and odd-order action corrections then explain how two later finite discrepancies can be repaired when their actual inputs and allowed changes are supplied.

\begin{longtable}{@{}>{\raggedright\arraybackslash}p{.23\textwidth}>{\raggedright\arraybackslash}p{.35\textwidth}>{\raggedright\arraybackslash}p{.36\textwidth}@{}}
\toprule
Statement & Status here & Boundary retained\\\midrule\endhead
Principal period and conductor & Ordinary derivation with [BH99] input & Does not construct a stratum\\[3pt]
Source self-pair $232$ & Expanded inherited ramification argument & Original source independence is retained\\[3pt]
Critical-four filter & Conditional actual-order consequence & Passing is not type occurrence or H5\\[3pt]
Rank-eight profile separation & Audited inherited construction and exact grade calculations & Three invariants do not force the native profile\\[3pt]
Classical source candidate & Selected primary-interface audit & Detailed [HL11] reread unavailable in this block\\[3pt]
Fixed-center extension criterion & Complete finite proof & Full sections can impose additional pointwise relations\\[3pt]
Free quadratic embedding & Complete finite proof & Requires the actual host refinement as input\\[3pt]
Transvection correction & Complete finite proof and exact tests & May change a preassigned embedding\\[3pt]
Odd-action correction & Complete conditional finite proof & Does not construct arithmetic actions or honest phases\\[3pt]
Native rank-eight refinement & Explicit remaining arithmetic obligation & Its value is not inferred from capacity\\[3pt]
General realization and H5 & Open in this project & No replacement by profile matching\\\bottomrule
\end{longtable}

\subsection{A certificate for the next native-marking step}
For a fixed arithmetic candidate at rank eight, a useful next certificate must identify the actual finite central extension or specified finite subgroup in its scalar pushout. It must give the chosen lift of each of the six symplectic basis vectors and compute its square under the full central character. These six bits determine the refinement once the commutator basis is fixed. Their three pairwise products determine the Arf sign through \eqref{eq:gauss}.

The certificate must also say whether the target marking is independently fixed or is allowed to be transported from the source. If it is fixed, compute the defect of the required quotient map. A vanishing defect directly certifies that map. An isotropic nonzero defect supplies the transvection only when changing the map is allowed. Every additional section element and transport lift then needs its actual multiplication check.

This task cannot be replaced by evaluating a canonical-looking quadratic form on an abstract six-dimensional space. That would choose new data rather than calculate the data of the specified arithmetic object. It also cannot be replaced by choosing another candidate with the same conductors. The candidate and full inducing character must remain fixed throughout.

\subsection{Verification and reproducibility boundaries}
The new standard-library exact suite contains 23 test methods. Its conductor tests cover the normalization table, grade dimensions, period divisibility, the rank-nine formal-degree identities, the rank-eight profiles and the finite field used for the projective source characters. Its marking tests cover extension group laws, square and section controls, Gauss signs, all labeled symplectic bases in dimension four, all linear-defect corrections through dimension six, relative and equivariant corrections, and odd versus even action cocycles.

The inherited 15-method Block 03 finite suite and 12-method Block 01 suite are replayed as separate regression gates. Their successful execution is not a new run of every historical arithmetic program. The supplied replay record lists exactly which commands were executed. Finite enumeration counts are not converted into a percentage of conjecture completion.

The source archive preserves the earlier compact checkpoint and versioned ledger history. It does not repair the missing Block 02 editable archive by inventing its contents. The v0.45 terminal compilation receipt remains unrecovered. No fresh Lean or Lake build occurred in this block.

The dedicated conjecture visual companion is the next publication task. Its mathematical route is now clear: normalized conductors, rejected and surviving rank-nine candidates, equal invariants with unequal capacity, transported versus native markings, the equal-capacity Arf gate, and the two conditional finite corrections. The original XI monograph, visual reader and Lean guide can then be synchronized against this stable claim map without presenting the general conjecture as a theorem.

\needspace{10\baselineskip}
\subsection{Source register}
The package's machine-readable source map supplies paths and hashes for inherited records. Page numbers below are printed pages unless a PDF-page locator is explicitly stated. Web-viewed primary sources are recorded as viewed sources, not as locally downloaded byte-pinned copies.

\begin{description}[style=nextline,leftmargin=1.5cm,labelwidth=1.35cm,font=\normalfont\bfseries,before=\raggedright]
\item[{[Q13]}] \emph{Adjoint Conductor Test161}. Retained dependency of Block166. The source invariant and conditional period-nine actual-order certificate. SHA-256 and exact source path are in the source map.
\item[{[Q14]}] \emph{Principal Order Conductor161}. Principal grades, leading-level normalization and necessary source matching.
\item[{[Q15]}] \emph{Arithmetic Constructor Audit}. The original source-defined induction and tame-descent branch, with full-character and source-input boundaries.
\item[{[Q16]}] \emph{Dyadic Markings167}. Full central extensions, transported finite markings, fixed native refinements and the sharp capacity criterion.
\item[{[D167]}] \emph{Dyadic Capacity167}. Full and skipped defining profiles, floor/ceiling grade counts, and the rank-eight shared-invariant example.
\item[{[BK]}] Colin J. Bushnell and Philip C. Kutzko. \emph{The Admissible Dual of GL(N) via Compact Open Subgroups}. Annals of Mathematics Studies 129, Princeton University Press, 1993. Relevant inherited inputs: recursive character orders, minimal commutator control, refinement, 6.2.5 and compact-type formal degree. The earlier selected audit remains distinguished from this block's fresh reading.
\item[{[BH99]}] Colin J. Bushnell and Guy Henniart. \emph{Local tame lifting for GL(n) II: wildly ramified supercuspidals}. Ast\'erisque 254, 1999. Freshly checked: introduction p. 2, 1.3 p. 15, 1.7 p. 18, 2.3 pp. 23--25, 6.1 pp. 53--54. The PDF page index is the printed page plus four when counted from zero.
\item[{[AP]}] Anne-Marie Aubert and Roger Plymen. \emph{Plancherel measure for GL(n,F) and GL(m,D): explicit formulas and Bernstein decomposition}. arXiv:math/0302169v4, 4 October 2004. Used: formulas (17), (29), Theorems 6.5 and 6.9, Remark 6.7 and the conductor convention in Section 4. The formal-degree pages inspected are pp. 19, 26--28.
\item[{[BH13]}] Colin J. Bushnell and Guy Henniart. \emph{To an effective local Langlands correspondence}. arXiv:1103.5316v2, 9 October 2013. The characteristic note on p. 8 is the selected input. Its effective construction is not used to supply its own totally wild starting correspondence.
\item[{[HL11]}] Guy Henniart and Bertrand Lemaire. \emph{Changement de base et induction automorphe pour GL(n) en caract\'eristique non nulle}. M\'emoires de la Soci\'et\'e Math\'ematique de France 124, 2011. DOI: 10.24033/msmf.436. The bibliographic record was checked. Detailed statements retain inherited audit status because the full PDF could not be retrieved in this block.
\end{description}
